% Part: first-order-logic
% Chapter: tableaux
% Section: derivations

\documentclass[../../../include/open-logic-section]{subfiles}

\begin{document}

\iftag{FOL}
      {\olfileid{fol}{tab}{der}}
      {\olfileid{pl}{tab}{der}}

\olsection{\usetoken{P}{tableau}}

\begin{explain}
We hebben uitgelegd wat een aanname is en kennen nu de regels waarmee
we formules afleiden. Daarmee bouwen we !!{tableau}s stap voor stap op.
Elk !!{tableau} ontstaat op een van twee manieren. Het is één tak met
een of meer aannamen. Of we maken het uit !!a{tableau} door op een tak
een van de afleidingsregels toe te passen.
\end{explain}

\begin{defn}[!!^{tableau}]
!!^a{tableau} voor de aannamen $\sFmla{S_1}{!A_1}$, \dots,
$\sFmla{S_n}{!A_n}$ is een eindige boom van !!{signed formula}s. Elke
$S_i$ is daarbij $\True$ of~$\False$. De boom moet aan de volgende
voorwaarden voldoen:
\begin{enumerate}
\item Bovenaan staan de $n$ !!{signed formula}s
  $\sFmla{S_i}{!A_i}$, onder elkaar.
\item Iedere andere !!{signed formula} in de boom moet ontstaan door
  een afleidingsregel correct toe te passen op !!a{signed formula} in
  de tak erboven.
\end{enumerate}
Een tak van !!a{tableau} noemen we precies dan \emph{gesloten} als hij
zowel $\sFmla{\True}{!A}$ als~$\sFmla{\False}{!A}$ bevat. Anders is de
tak \emph{open}. !!^a{tableau} waarin elke tak gesloten is, noemen we
een \emph{gesloten !!{tableau}} voor de gebruikte aannamen. Als
!!a{tableau} niet gesloten is, bevat het minstens één open tak. Dan is
het tableau \emph{open}.
\end{defn}

\begin{ex}
  Een verzameling aannamen is op zichzelf al !!a{tableau}. Meestal is
  dat tableau nog niet gesloten. Het is alleen gesloten als de aannamen
  al het volgende paar !!{signed formula}s bevatten:
  $\sFmla{\True}{!A}$ en~$\sFmla{\False}{!A}$.

  Uit !!a{tableau}, open of gesloten, kunnen we een nieuw en groter
  tableau maken. Kies daarin !!a{signed formula}~$!A$ en pas er een
  afleidingsregel op toe. De regel zet een of meer nieuwe
  !!{signed formula}s onderaan iedere tak die precies dat voorkomen
  van~$!A$ bevat.

  Neem bijvoorbeeld de aanname $\sFmla{\True}{!A \land \lnot
    !A}$. Eerst hebben we alleen die aanname. Dat geeft het volgende
  open !!{tableau}:
  \begin{center}
    \begin{tableau}{}
      [\sFmla{\True}{\formula{A} \land \lnot \formula{A}}, just=\TAss]
    \end{tableau}{}
  \end{center}
  Nu passen we de regel $\TRule{\True}{\land}$ op de aanname toe. Met
  die regel mogen we twee nieuwe regels aan het !!{tableau} toevoegen:
  $\sFmla{\True}{!A}$ en $\sFmla{\True}{\lnot !A}$.
  \begin{center}
    \begin{tableau}{}
      [\sFmla{\True}{\formula{A} \land \lnot \formula{A}}, just=\TAss
        [\sFmla{\True}{\formula{A}}, just={\TRule{\True}{\land}[1]},
          [\sFmla{\True}{\lnot \formula{A}}, just={\TRule{\True}{\land}[1]}
          ]
        ]
      ]
    \end{tableau}{}
  \end{center}
  Rechts schrijven we bij iedere regel van !!a{tableau} welke
  afleidingsregel we hebben gebruikt. Zo betekent
  $\TRule{\True}{\land} 1$ dat de !!{signed formula} op die regel is
  ontstaan door $\TRule{\True}{\land}$ toe te passen op de
  !!{signed formula} op regel~$1$. Het nieuwe !!{tableau} bevat nu meer
  !!{signed formula}s. Op slechts één daarvan,
  $\sFmla{\True}{\lnot !A}$, kunnen we nog een regel toepassen. Dat is
  hier de regel $\TRule{\True}{\lnot}$. Na toepassing daarvan is het
  !!{tableau} gesloten:
  \begin{center}
    \begin{tableau}{}
      [\sFmla{\True}{\formula{A} \land \lnot \formula{A}}, just=\TAss
        [\sFmla{\True}{\formula{A}}, just={\TRule{\True}{\land}[1]}
          [\sFmla{\True}{\lnot \formula{A}}, just={\TRule{\True}{\land}[1]}
            [\sFmla{\False}{\formula{A}}, just={\TRule{\True}{\lnot}[3]}, close]
          ]
        ]
      ]
    \end{tableau}{}
  \end{center}
\end{ex}

\end{document}
